\documentclass[11pt]{article}
\usepackage{amsmath,amssymb,amsthm}
\usepackage[margin=0.8in]{geometry}
\usepackage[colorlinks=true,linkcolor=blue,urlcolor=blue]{hyperref}

\newtheorem{theorem}{Theorem}
\newtheorem{lemma}{Lemma}
\newtheorem{proposition}{Proposition}
\newtheorem{corollary}{Corollary}
\newtheorem{remark}{Remark}

\newcommand{\Z}{\mathbb{Z}}
\newcommand{\N}{\mathbb{N}}

\title{Prescribed Isolated Intervals in Higher Fold Sumsets}
\author{Zijie Yu\\
Ningbo University\\
Ningbo, China\\
\texttt{256001428@nbu.edu.cn}}
\date{June 5, 2026}

\begin{document}
\maketitle

\begin{abstract}
We give an explicit construction of a finite set of integers whose $h$-fold
sumset has any prescribed sufficiently separated family of nontrivial
intervals and no others. This answers Problems 9 and 10 in
\href{https://arxiv.org/abs/2605.26425}{Nathanson's problem list} with an
explicit set of size $2q(n+1)$, and also answers the higher-order Problem 11.
We then replace the initial exponential labels by finite moment labels. For
fixed $h$, this gives polynomial diameter in the number of prescribed target
points and works for an arbitrary finite target set. Both label choices are
instances of a common separated-label principle. For double sumsets, a
Sidon-label instance improves the diameter bound to quadratic. We record a
small counterexample to the strict inequality proposed in Problem 7(2).
\end{abstract}

\section{A short separation lemma}

For an integer vector $v=(v_0,\ldots,v_{s-1})$, write
$\lVert v\rVert_1=\sum_r |v_r|$.

\begin{lemma}[Positional separation]\label{lem:separation}
Let $B,M\geq 1$ and $a_r=MB^r$ for $0\leq r<s$. If
$\lambda,\mu\in\Z^s$, $\lambda\neq\mu$, and
$\lVert\lambda\rVert_1+\lVert\mu\rVert_1\leq B$, then
\[
 \left|\sum_{r=0}^{s-1}(\lambda_r-\mu_r)a_r\right|\geq M.
\]
\end{lemma}

\begin{proof}
Put $z=\lambda-\mu$ and let $R$ be the largest index with $z_R\neq 0$.
We have $\lVert z\rVert_1\leq B$. If $R=0$, the result is immediate. If
$R\geq 1$, then
\begin{align*}
\left|\sum_{r=0}^{s-1}z_ra_r\right|
&\geq M\left(|z_R|B^R-\sum_{r<R}|z_r|B^r\right)\\
&\geq M\left(B^R-(B-1)B^{R-1}\right)
\geq M.
\end{align*}
Here the sum of the absolute values of the lower coefficients is at most
$B-|z_R|\leq B-1$.
\end{proof}

\begin{proposition}[Separated-label principle]\label{prop:labels}
Let $h\geq 2$, $d\geq 1$, and let
$S=\{t_0,\ldots,t_{m-1}\}\subseteq\Z$ be nonempty. Put
\[
 H=h(h-1),
 \qquad
 T=\max_{t\in S}|t|.
\]
Suppose that $M\geq 2hT+d$ and that integers $a_0,\ldots,a_{m-1}$ satisfy
\[
 \left|\sum_{r=0}^{m-1}(\lambda_r-\mu_r)a_r\right|\geq M
 \tag{1}\label{eq:abstract-separation}
\]
whenever $\lambda,\mu\in\Z^m$ are distinct and
\[
 \lVert\lambda\rVert_1+\lVert\mu\rVert_1\leq 2H.
\]
Then
\[
 A=\{t_r+(h-1)a_r:0\leq r<m\}\cup\{-a_r:0\leq r<m\}
\]
has $|A|=2m$ and
\[
 hA=S\mathbin{\dot\cup}E,
\]
where distinct elements of $E$ are at distance at least $d$ from one another
and every element of $E$ is at distance at least $d$ from $S$.
\end{proposition}

\begin{proof}
Attach to $t_r+(h-1)a_r$ the label vector $(h-1)e_r$ and offset $t_r$, and
attach to $-a_r$ the label vector $-e_r$ and offset $0$. Distinct tagged
generators have distinct label vectors, combined norm at most $2H$, and
offsets differing by at most $2T<M$. Thus \eqref{eq:abstract-separation}
shows that their integer values are distinct, so $|A|=2m$.

Every element of $hA$ is a bounded offset plus $\sum_r v_ra_r$, where
$\lVert v\rVert_1\leq H$. Let $p_r$ and $u_r$ count the positive and negative
terms with index $r$ in an $h$-term sum. A zero label vector satisfies
$(h-1)p_r=u_r$ for every $r$. Since
\[
 \sum_r(p_r+u_r)=h\sum_rp_r=h,
\]
there is exactly one positive term and $h-1$ matching negative terms. Hence
the zero-label sums are exactly
\[
 (t_r+(h-1)a_r)+(h-1)(-a_r)=t_r.
\]

Every other sum has a nonzero label vector. Such a vector determines its
multiset of terms: if two multisets with counts $(p_r,u_r)$ and
$(p'_r,u'_r)$ had the same nonzero vector but were different, then
\[
 u_r-u'_r=(h-1)(p_r-p'_r).
\]
Summing and using that both multisets have $h$ terms shows that
$\sum_rp_r=\sum_rp'_r$. At an index with $p_r>p'_r$, the first multiset
would contain at least one positive term and $h-1$ matching negative terms.
It would therefore be the zero-label matched sum, a contradiction.

Finally, every representation offset has absolute value at most $hT$.
Distinct exceptional sums have distinct label vectors, and an exceptional
sum also has a label vector distinct from the zero vector of a target.
Equation~\eqref{eq:abstract-separation} and the offset bound leave a gap of
at least $M-2hT\geq d$.
\end{proof}

\section{The construction}

\begin{theorem}[Prescribed isolated intervals]\label{thm:main}
Let $h\geq 2$, $n\geq 1$, $d\geq 2$, $q\geq 1$, and
\[
 C=\{c_1<\cdots<c_q\}\subseteq\Z
\]
with $c_{i+1}-c_i\geq n+2$ for all $i$. There is an explicitly constructible
set $A\subseteq\Z$ with $|A|=2q(n+1)$ such that
\[
 hA=S\mathbin{\dot\cup}E,
 \qquad
 S=\bigcup_{i=1}^q[c_i,c_i+n],
\]
where
\[
 |x-y|\geq d
\]
whenever $x,y\in E$ are distinct or $x\in E$ and $y\in S$.
Consequently, the maximal intervals in $hA$ of length at least one are
exactly the intervals $[c_i,c_i+n]$.
\end{theorem}

\begin{proof}
List the $q(n+1)$ elements of $S$ as $t_0,\ldots,t_{m-1}$, where
$m=q(n+1)$. Set
\[
 T=\max_{t\in S}|t|,
 \qquad M=2hT+d,
 \qquad B=2h(h-1),
 \qquad a_r=MB^r,
\]
and define
\[
 A=\{t_r+(h-1)a_r:0\leq r<m\}\cup\{-a_r:0\leq r<m\}.
\]
Lemma~\ref{lem:separation} verifies \eqref{eq:abstract-separation} with
$B=2h(h-1)$. Proposition~\ref{prop:labels} gives $|A|=2m$, the decomposition
$hA=S\mathbin{\dot\cup}E$, and the required distance bounds.
Since $d\geq 2$ and the intervals in $S$ are separated by at least one
missing integer, no exceptional point can extend a target interval or form a
new nontrivial interval.
\end{proof}

\begin{corollary}[Single interval replacement]\label{cor:single}
For all $h\geq 2$, $c\in\Z$, $n\geq 1$, and $d\geq 2$, there is a set
$A\subseteq\Z$ with $|A|=2n+2$ such that $[c,c+n]\subseteq hA$, every
element of $hA\setminus[c,c+n]$ is at distance at least $d$ from every
other such element and from $[c,c+n]$, and the longest interval in $hA$ has
length $n$.
\end{corollary}

\begin{proof}
Apply Theorem~\ref{thm:main} with $C=\{c\}$.
\end{proof}

\begin{corollary}[Problems 9--11]\label{cor:problems}
In the notation of Problems 9--11 of Nathanson's problem list, the
assertions hold for every $N_q$ allowed in the question, with
\[
 k_C=2q(n+1).
\]
\end{corollary}

\begin{proof}
The hypotheses give $c_{i+1}-c_i\geq N_q\geq n+2$. Apply
Theorem~\ref{thm:main} with $d=2$, taking $h=2$ for Problems 9--10 and the
specified value of $h$ for Problem 11.
\end{proof}

\begin{remark}[Diameter bound]
The construction in Theorem~\ref{thm:main} satisfies
\[
 \operatorname{diam}(A)
 \leq h(2hT+d)\bigl(2h(h-1)\bigr)^{q(n+1)-1}+T.
\]
The bound is intentionally simple rather than optimized.
\end{remark}

\section{A polynomial diameter refinement}

The power labels in Theorem~\ref{thm:main} make the first construction
transparent, but their size is exponential in the number of target points.
The following finite moment encoding gives a polynomial bound for fixed $h$.

\begin{lemma}[Moment separation]\label{lem:moment-separation}
Let $m,L\geq 1$, put
\[
 R=\max(1,m-1),
 \qquad
 K=LR^{L-1},
 \qquad
 Q=LK=L^2R^{L-1},
\]
and define
\[
 b_r=\sum_{j=0}^{L-1}r^jQ^j
 \qquad (0\leq r<m).
\]
If $z=(z_0,\ldots,z_{m-1})\in\Z^m$ is nonzero and
$\lVert z\rVert_1\leq L$, then
\[
 \left|\sum_{r=0}^{m-1}z_rb_r\right|\geq 1.
\]
\end{lemma}

\begin{proof}
For $0\leq j<L$, put
\[
 \mu_j=\sum_{r=0}^{m-1}z_rr^j.
\]
Not all $\mu_j$ vanish. Indeed, if the support of $z$ has size $s$, then
$s\leq L$, and the first $s$ equations $\mu_j=0$ would form an invertible
Vandermonde system on the distinct support indices.

Since $|\mu_j|\leq LR^{L-1}=K$, the moment vector has
$\ell^1$ norm at most $LK=Q$. A nonzero integer coefficient vector with
$\ell^1$ norm at most $Q$ cannot evaluate to zero in base $Q$: if $J$ is
its largest nonzero index, the leading term has absolute value at least
$Q^J$, while all lower terms together have absolute value at most
$(Q-1)Q^{J-1}$ when $J\geq 1$; the case $J=0$ is immediate. Applying this
positional domination to $(\mu_0,\ldots,\mu_{L-1})$ proves the claim.
\end{proof}

\begin{theorem}[Polynomial diameter realization]\label{thm:polynomial}
Let $h\geq 2$, $d\geq 2$, and let $S\subseteq\Z$ be a nonempty finite set
with $m=|S|$. There is an explicitly constructible set $A\subseteq\Z$ with
$|A|=2m$ such that
\[
 hA=S\mathbin{\dot\cup}E,
\]
where
\[
 |x-y|\geq d
\]
whenever $x,y\in E$ are distinct or $x\in E$ and $y\in S$.

More precisely, let $\tau\in h\Z$ minimize
\[
 T=\max_{t\in S}|t-\tau|.
\]
Put
\[
 L=2h(h-1),
 \quad
 R=\max(1,m-1),
 \quad
 Q=L^2R^{L-1},
\]
and
\[
 P_h(m)=\sum_{j=0}^{L-1}(m-1)^jQ^j.
\]
Then $A$ can be chosen with
\[
 \operatorname{diam}(A)
 \leq h(2hT+d)P_h(m)+T.
\]
For fixed $h$, the factor $P_h(m)$ is bounded by a polynomial in $m$ of
degree $L(L-1)$.
\end{theorem}

\begin{proof}
Enumerate $S=\{t_0,\ldots,t_{m-1}\}$ and use the parameters in
Lemma~\ref{lem:moment-separation}. Set
\[
 M=2hT+d,
 \qquad
 a_r=M\sum_{j=0}^{L-1}r^jQ^j,
\]
and first construct the centered set
\[
 A_0=
 \{t_r-\tau+(h-1)a_r:0\leq r<m\}
 \cup
 \{-a_r:0\leq r<m\}.
\]
Finally, put
\[
 A=A_0+\tau/h.
\]

If $\lambda\neq\mu$ and
$\lVert\lambda\rVert_1+\lVert\mu\rVert_1\leq L$, then
$z=\lambda-\mu$ is nonzero and $\lVert z\rVert_1\leq L$.
Lemma~\ref{lem:moment-separation}, after multiplication by $M$, therefore
verifies \eqref{eq:abstract-separation} for the centered target set.
Proposition~\ref{prop:labels} gives $|A_0|=2m$, the matching-label
decomposition, and the required distance bounds.
Translating $A_0$ by $\tau/h$ translates $hA_0$ by $\tau$ and does not change
its diameter.

The labels $a_r$ are strictly increasing. Every element of $A_0$ lies between
$-a_{m-1}$ and $T+(h-1)a_{m-1}$. Consequently,
\[
 \operatorname{diam}(A)=\operatorname{diam}(A_0)
 \leq ha_{m-1}+T
 =h(2hT+d)P_h(m)+T.
\]
For $m\geq 2$, the displayed expression for $P_h(m)$ is a polynomial of
degree $L(L-1)$; the case $m=1$ is immediate.
\end{proof}

\begin{corollary}[Translation-invariant interval bound]\label{cor:poly-interval}
Under the hypotheses of Theorem~\ref{thm:main}, let
\[
 m=q(n+1),
 \qquad
 D=c_q+n-c_1.
\]
There is a set $A$ with the same isolated interval pattern and
$|A|=2q(n+1)$ whose diameter is polynomial in $m$ for fixed $h$, with
\[
 T\leq \left\lceil\frac{D+h}{2}\right\rceil
\]
in the bound of Theorem~\ref{thm:polynomial}. In particular, the bound is
independent of the absolute location of the prescribed intervals.
\end{corollary}

\begin{proof}
Apply Theorem~\ref{thm:polynomial} to
$S=\bigcup_i[c_i,c_i+n]$. A multiple $\tau$ of $h$ can be chosen within
$h/2$ of the midpoint of $[c_1,c_q+n]$, giving the stated bound for $T$.
\end{proof}

\begin{theorem}[Quadratic diameter for double sumsets]\label{thm:sidon}
Let $d\geq 2$, and let $S\subseteq\Z$ be a nonempty finite set with $m=|S|$.
There is an explicitly constructible set $A\subseteq\Z$ with $|A|=2m$ such
that
\[
 2A=S\mathbin{\dot\cup}E,
\]
where distinct elements of $E$ are at distance at least $d$ from one another
and every element of $E$ is at distance at least $d$ from $S$.

More precisely, let $\tau\in 2\Z$ minimize
\[
 T=\max_{t\in S}|t-\tau|,
\]
and let $p\geq m$ be an odd prime. For $0\leq r<m$, let $\rho_r$ be the least
nonnegative residue of $r^2$ modulo $p$, and put
\[
 b_r=2p^2+2pr+\rho_r.
\]
Then $A$ can be chosen with
\[
 \operatorname{diam}(A)\leq 8(4T+d)p^2+T.
\]
For $m\geq 2$, Bertrand's postulate allows $p<2m$, giving a quadratic bound
in $m$. Since $p$ is an integer, the resulting explicit estimate is
\[
 \operatorname{diam}(A)
 \leq 8(4T+d)(2m-1)^2+T.
\]
\end{theorem}

\begin{proof}
First observe that $\{b_0,\ldots,b_{m-1}\}$ is a Sidon set, with repetitions
allowed in pair sums. Indeed, suppose
\[
 b_i+b_j=b_k+b_\ell.
\]
Reducing modulo $2p$ gives $\rho_i+\rho_j=\rho_k+\rho_\ell$, because both sides
lie between $0$ and $2p-2$. It follows that $i+j=k+\ell$. Reducing the residue
identity modulo $p$ also gives
\[
 i^2+j^2\equiv k^2+\ell^2\pmod p.
\]
Since $p$ is odd, we obtain $ij\equiv k\ell\pmod p$. Thus the two pairs are the
roots of the same quadratic polynomial over $\mathbb F_p$. The unordered
pairs $\{i,j\}$ and $\{k,\ell\}$ agree modulo $p$, hence agree as integer pairs
because all indices lie in $[0,p)$.

Moreover,
\[
 2p^2\leq b_r<4p^2.
\]
Consequently, an equality between two sums of the $b_r$ involving at most four
terms in total must have the same number of terms on both sides: one term is
strictly smaller than any sum of two terms. Equal one-term sums are trivial,
and equal two-term sums are trivial by the Sidon property. Therefore
\[
 \sum_r z_rb_r\neq 0
\]
for every nonzero integer vector $z$ with $\lVert z\rVert_1\leq 4$.

Set $M=4T+d$, enumerate $S=\{t_0,\ldots,t_{m-1}\}$, and apply
Proposition~\ref{prop:labels} to the centered targets $t_r-\tau$ and the labels
$a_r=Mb_r$. Finally translate the resulting basis by $\tau/2$. Since
$b_r<4p^2$, its diameter is at most
\[
 2M\max_r b_r+T\leq 8(4T+d)p^2+T.
\]
The same construction has quadratic diameter from below. Indeed, translation
does not change diameter, the negative generators contain
$-Mb_0$ and $-Mb_{m-1}$, and
\[
 b_{m-1}-b_0
 =2p(m-1)+\rho_{m-1}
 \geq 2m(m-1).
\]
Consequently,
\[
 \operatorname{diam}(A)\geq 2M m(m-1).
\]
\end{proof}

\begin{proposition}[A Sidon-label diameter barrier]\label{prop:sidon-barrier}
Let $B=\{b_1,\ldots,b_m\}\subseteq\Z$ be a Sidon set, where repetitions are
allowed in pair sums. Then
\[
 \operatorname{diam}(B)\geq \binom{m}{2}.
\]
Consequently, any tagged construction that contains the generators
$-Mb_i$ has diameter at least $M\binom{m}{2}$.
\end{proposition}

\begin{proof}
Order the labels so that $b_1<\cdots<b_m$. The $\binom{m}{2}$ positive
differences $b_j-b_i$, with $i<j$, are distinct. Indeed, an equality
$b_j-b_i=b_\ell-b_k$ gives $b_j+b_k=b_\ell+b_i$, and the Sidon property
forces the two ordered index pairs to agree. Every positive difference is at
most $\operatorname{diam}(B)$, proving the first claim. The second follows
because the generators $-Mb_i$ already have diameter
$M\operatorname{diam}(B)$.
\end{proof}

\begin{remark}[Elementary lower bounds]\label{rem:lower-bounds}
Let $A\subseteq\Z$ be finite and nonempty, put $k=|A|$, and suppose that
$S\subseteq hA$ has $m$ elements. Then
\[
 m\leq \binom{k+h-1}{h},
 \qquad
 \operatorname{diam}(A)\geq \frac{\operatorname{diam}(S)}{h}.
\]
Indeed, the first inequality counts unordered $h$-term multisets from $A$,
and the second follows from
\[
 \operatorname{diam}(hA)=h\operatorname{diam}(A).
\]
Thus Theorem~\ref{thm:polynomial} supplies a polynomial diameter bound for
fixed $h$, while Theorem~\ref{thm:sidon} gives a quadratic bound when $h=2$.
Proposition~\ref{prop:sidon-barrier} shows that the quadratic order is
unavoidable within the Sidon-label method. Optimizing either the cardinality
or the degree for $h\geq 3$ remains open.
\end{remark}

\section{Audit notes and a small counterexample}

The proof above also supplies a replacement proof for the single-interval
statement in Theorem 4 of Nathanson's preprint. In the current source, the
displayed separation parameter is
\[
 \Delta=\min(3,2|c|+d+2n),
\]
while later estimates require a parameter larger than quantities depending
on $|c|$, $d$, and $n$. The replacement argument avoids that discrepancy.

Problem 7(2) asks whether a negative element forces a strict loss in the
initial interval. This is false: take $h=2$, $k=3$, and
\[
 A=\{-1,1,2\}.
\]
Then $[0,4]\subseteq 2A$. To see that $n_2(3)=4$, write a nonnegative
three-element basis for an initial interval as $\{0,1,a\}$. If $a=2$ or
$a=3$, its initial interval ends at $4$, while if $a>3$, the integer $3$ is
missing. Hence $\ell_2(A)=n_2(3)$ even though $\min(A)<0$.

\section{Reproducibility}

The accompanying Python package constructs $A$, computes $hA$, checks all
maximal consecutive runs, records the minimum exceptional spacing, and emits
JSON certificates. It also includes exact small-parameter computations
cross-checked against \href{https://oeis.org/A001208}{OEIS A001208} and
\href{https://oeis.org/A001209}{OEIS A001209}. The arithmetic separation,
bounded-offset, matched-count, tagged-representation, exceptional-count
uniqueness, generic set-level isolation, interval-enumeration, maximal-run,
and exact-cardinality arguments have also been checked in Lean 4 with Mathlib.
The resulting Lean theorem automatically selects the target radius and
substitutes the displayed values of $M$ and $B$. The finite moment-label
Vandermonde argument, translation rule, closed diameter estimate, exact
cardinality, midpoint-centering lemma, and interval-family theorem have also
been formalized. The audit note about the earlier preprint remains an
ordinary paper argument. For the shifted-Sidon refinement, Lean checks the
size bounds and the reduction from
order-four short-relation freedom to set-level isolation, including the
explicit pair-sum argument over $\mathbb F_p$. It also checks translation
invariance, quadratic diameter upper and lower bounds, tightened
Bertrand-prime selection, exact cardinality after translation, and the final
separated-interval wrapper. The midpoint translation, radius, and scale are
selected automatically from an enclosing interval. Lean also checks the direct
Problem 7(2)
witness $\{-1,1,2\}$ and the containment $[0,4]\subseteq 2A$; the classical
identification $n_2(3)=4$ remains the short paper argument given above.

\end{document}
